\documentclass[11pt,letterpaper]{article}
\usepackage[utf8]{inputenc}
\usepackage[margin=1in]{geometry}
\usepackage{booktabs}
\usepackage{longtable}
\usepackage{hyperref}
\usepackage{listings}
\usepackage{xcolor}
\usepackage{enumitem}

\definecolor{codegreen}{rgb}{0,0.6,0}
\definecolor{codegray}{rgb}{0.5,0.5,0.5}
\definecolor{backcolour}{rgb}{0.95,0.95,0.92}
\definecolor{cautionred}{rgb}{0.55,0,0}
\definecolor{okgreen}{rgb}{0,0.45,0}

\lstset{
    backgroundcolor=\color{backcolour},
    commentstyle=\color{codegreen},
    basicstyle=\ttfamily\footnotesize,
    breaklines=true, captionpos=b, keepspaces=true,
    numbers=none, showspaces=false, showstringspaces=false, tabsize=2
}

\newenvironment{caution}[1]{%
    \par\addvspace{8pt}%
    \noindent{\color{cautionred}\rule{\linewidth}{0.8pt}}\par\vspace{2pt}
    \noindent\textbf{\color{cautionred}Caution: #1}\par\vspace{2pt}
}{%
    \par\vspace{2pt}\noindent{\color{cautionred}\rule{\linewidth}{0.4pt}}%
    \par\addvspace{8pt}%
}

\newcommand{\passif}[1]{\par\vspace{2pt}\noindent\textbf{\color{okgreen}Pass:} #1\par\vspace{4pt}}
\newcommand{\failif}[1]{\noindent\textbf{If it fails:} #1\par\vspace{6pt}}

\title{SOP: Bringing Up a NEMA 17 Joint \\ \large Arctos Robotic Arm --- MKS Servo42D over CAN}
\author{Applies to Joints A, B and C}
\date{September 2026}

\begin{document}
\maketitle

\noindent This is the step-by-step procedure for taking a NEMA 17 joint from
unpowered to ready for use. It is procedure only --- for protocol details, encoder
maths, wiring and troubleshooting, see \texttt{arctos\_motor\_guide.tex}.

\medskip
\noindent\textbf{Do the steps in order.} Each one gates the next: the procedure moves
from information to motion, and never commands motion before the joint has been
proven to answer correctly.

\begin{table}[h]
\centering
\begin{tabular}{llll}
\toprule
\textbf{Joint} & \textbf{CAN ID} & \textbf{Gear ratio} & \textbf{Role} \\ \midrule
A & \texttt{0x04} & 48.0:1  & Elbow \\
B & \texttt{0x05} & 67.82:1 & Wrist rotation \\
C & \texttt{0x06} & 67.82:1 & Wrist joint \\
\bottomrule
\end{tabular}
\caption{Substitute your joint's values wherever the procedure says
\texttt{<ID>} and \texttt{<RATIO>}. Examples below use Joint A.}
\end{table}

\section*{Step 0 --- Before applying power}
\begin{enumerate}[leftmargin=*]
    \item Confirm the controller is the \textbf{CAN variant} of the MKS Servo42D.
          Read the part number on the board. An RS485 variant has no CAN transceiver
          and will never respond, no matter what you do in software.
    \item With everything powered off, measure \texttt{CAN\_H} to \texttt{CAN\_L}.
          Expect $\approx 60\,\Omega$.
    \item Confirm only \texttt{CAN\_H}, \texttt{CAN\_L} and \texttt{GND} run between
          the CANable and the controller.
    \item Confirm the joint is mechanically clear and, if it is on the arm, that you
          know its travel limits.
\end{enumerate}
\passif{$\approx 60\,\Omega$ across the bus and a confirmed CAN-variant board.}
\failif{$120\,\Omega$ means one terminator; open means broken wiring. Fix before continuing.}

\section*{Step 1 --- Configure the controller panel}
Power the controller. On its display, set and \textbf{save}:

\begin{table}[h]
\centering
\begin{tabular}{ll}
\toprule
CAN ID & \texttt{<ID>} for this joint \\
CAN Rate & \texttt{500} \\
CAN Response (CANRSP) & \texttt{ENABLE} \\
CAN CRC & \texttt{SUM-8} \\
Working current & $\leq 1.2$~A (the NEMA 17's rating) \\
\bottomrule
\end{tabular}
\end{table}

Power-cycle the controller after saving.

\begin{caution}{Check the working current}
Panel defaults have been found as high as 3000~mA against a 1.2~A motor. Set this
explicitly.
\end{caution}
\passif{Settings persist across the power cycle.}

\section*{Step 2 --- Bring up the CAN interface}
\begin{lstlisting}[language=bash]
lsusb                                    # expect ID 16d0:117e CANable2
ls -l /dev/serial/by-id/                 # note the stable device name
sudo pkill slcand
sudo ip link delete can0 2>/dev/null
sudo slcand -o -c -s6 /dev/serial/by-id/<your-canable> can0
sudo ip link set can0 txqueuelen 1000
sudo ip link set up can0
ip -br link show can0
\end{lstlisting}
\passif{\texttt{can0 \quad UP}}
\failif{Re-check the device name --- it changes whenever the adapter re-enumerates.
\texttt{bitrate 0} in \texttt{ip -details} output is normal and not a fault.}

\section*{Step 3 --- Confirm the real CAN ID}
Never assume the ID from a config file or a table.
\begin{lstlisting}[language=bash]
python3 src/arctos_can_control/arctos_can_control/can_id_scan.py \
    --channel can0 --ids 1-16
\end{lstlisting}
This sends only the read-only \texttt{0x31} query. It never enables coils and never
commands motion, so it is safe at any time.
\passif{Exactly one responder, and its ID is the one you set in Step 1.}
\failif{No responders: re-check panel settings, wiring, and the board's part number.
Multiple responders: another driver is powered on the bus --- resolve the collision
before going further.}

\section*{Step 4 --- Read-only health check}
With the ID confirmed, verify the link is trustworthy before commanding anything.
\begin{lstlisting}[language=bash]
python3 src/arctos_can_control/arctos_can_control/joint_reliability_test.py \
    --channel can0 --can-id <ID> --gear-ratio <RATIO> --samples 200
\end{lstlisting}
Check the following, all with the joint at rest:

\begin{table}[h]
\centering
\begin{tabular}{lll}
\toprule
\textbf{Query} & \textbf{Opcode} & \textbf{Expected} \\ \midrule
Encoder        & \texttt{0x31} & stable; zero jitter, zero drift \\
Velocity       & \texttt{0x32} & zero \\
Error          & \texttt{0x39} & small (tens of counts) \\
Enable state   & \texttt{0x3A} & \texttt{01} --- these drivers power up enabled \\
Shaft protection & \texttt{0x3E} & \texttt{00} \\
\bottomrule
\end{tabular}
\end{table}
\passif{100\% response rate, 0 checksum failures, 0 encoder jitter at rest.}
\failif{Any checksum failure or dropped response is a wiring or bitrate problem.
Do not proceed to motion.}

\section*{Step 5 --- Check the mechanism by hand}
Do this before the motor ever drives the joint. It separates mechanical faults from
electrical ones, which the bus alone cannot distinguish.

\begin{enumerate}[leftmargin=*]
    \item Disable the coils: \texttt{cansend can0 <ID>\#F300<crc>}.
    \item Confirm with \texttt{0x3A} that the state is now \texttt{00}.
    \item Turn the joint output by hand, gently, both directions, while watching the
          encoder.
\end{enumerate}

\begin{caution}{Only drop coils if the joint is not gravity-loaded}
On an assembled arm this removes holding torque and the joint may fall. Support it,
or skip this step if it cannot be done safely.
\end{caution}

\passif{The encoder tracks your hand in both directions, and the motion feels smooth.
This proves the gearbox transmits and nothing is skipping.}
\failif{Gritty or catching resistance means the gear mesh needs lubrication --- use a
plastic-safe PTFE or silicone grease. No encoder movement while the output turns
means the motor shaft is not following: check the coupling, grub screw, or bore.}

\section*{Step 6 --- First commanded motion}
Re-enable coils, then command one small move. \textbf{Have the emergency stop ready
in a second terminal before you send anything:}
\begin{lstlisting}[language=bash]
# Joint A example -- recompute the checksum for your joint
cansend can0 004#F7FB
\end{lstlisting}

\begin{lstlisting}[language=bash]
python3 src/arctos_can_control/arctos_can_control/joint_unstick_move_test.py \
    --can-id <ID> --gear-ratio <RATIO> \
    --step-degrees 0.25 --num-steps 1 --max-unstick 0 \
    --band-lo -1 --band-hi 1
\end{lstlisting}

\begin{caution}{Set the band explicitly}
Some joints have \texttt{home\_position} and \texttt{opposite\_limit} of \texttt{0.0}
in \texttt{arctos\_controller.yaml}, which produces an inverted, unusable safe band.
Always pass \texttt{--band-lo}/\texttt{--band-hi} yourself.
\end{caution}

\passif{The step reports $\approx$100\% of commanded and the status frames show
\texttt{FD=[1, 2]}. Both frames matter: \texttt{FD 01} alone means the command
started but did not finish.}
\failif{Do not retry blindly --- that only heats the motor. Go back to Step 5.}

\section*{Step 7 --- Stepped travel test}
Extend to real travel, one direction then the other. Use a step size and count that
stay inside the joint's safe range.
\begin{lstlisting}[language=bash]
# forward
python3 src/arctos_can_control/arctos_can_control/joint_unstick_move_test.py \
    --can-id <ID> --gear-ratio <RATIO> \
    --step-degrees 2.0 --num-steps 9 --max-unstick 0 \
    --band-lo -25 --band-hi 25 --speed 25

# reverse -- same step count, to close the loop
python3 src/arctos_can_control/arctos_can_control/joint_unstick_move_test.py \
    --can-id <ID> --gear-ratio <RATIO> \
    --step-degrees -2.0 --num-steps 9 --max-unstick 0 \
    --band-lo -25 --band-hi 25 --speed 25
\end{lstlisting}
\passif{Every step 99--101\% of commanded with \texttt{FD 02}, and the round trip
returns to the starting position. On a healthy joint this closes to within a few
encoder counts.}
\failif{A step that under-delivers and then a step that delivers nothing means the
joint is against something, or the motor has lost drive. Stop and diagnose --- do
not increase torque or retry.}

\section*{Step 8 --- Acceptance}
The joint is ready for use when all of the following hold:

\begin{table}[h]
\centering
\begin{tabular}{ll}
\toprule
\textbf{Check} & \textbf{Criterion} \\ \midrule
CAN ID           & confirmed by scan, matches config \\
Comms            & 100\% response, 0 checksum failures \\
Encoder at rest  & 0 jitter, 0 drift \\
Hand check       & smooth, encoder tracks both directions \\
Per-step accuracy & 99--101\% of commanded \\
Status frames    & \texttt{FD 02} on every step \\
Round trip       & returns to start within a few counts \\
Shaft protection & still \texttt{00} \\
\bottomrule
\end{tabular}
\end{table}

\section*{Step 9 --- Operating limits}
Once accepted, keep within these when using the joint:

\begin{itemize}[leftmargin=*]
    \item \textbf{Speed at or below 225} (motor rpm) --- \emph{for a bare motor}.
          Above $\approx$250 an unloaded motor receives more pulses than it
          executes and silently loses position; it saturates near 310~rpm. These
          are no-load figures and a geared joint saturates well below them: Joint
          B, through 67.82:1, tracks only to 150~rpm and plateaus near 165, though
          it did not lose position when over-commanded. Measure the joint rather
          than inheriting this number.
    \item \textbf{Use \texttt{0xFD} for all motion.} Do not use \texttt{0xF5}.
    \item \textbf{Never retry into a stall.} Back off and re-attempt, or stop.
    \item \textbf{Disable coils between tests} so the motor is not needlessly
          energized --- unless the joint is gravity-loaded, in which case leave it
          holding.
\end{itemize}

\section*{Step 10 --- Shutdown}
\begin{lstlisting}[language=bash]
# stop any running script first (Ctrl+C)
cansend can0 <ID>#F300<crc>     # disable coils, if safe to do so
sudo ip link set can0 down
sudo pkill slcand
# then unplug the adapter
\end{lstlisting}

\end{document}
